\documentclass[conference]{IEEEtran}
\renewcommand{\refname}{Referências}
\IEEEoverridecommandlockouts
\usepackage{cite}
\usepackage{amsmath,amssymb,amsfonts}
\usepackage{algorithmic}
\usepackage{graphicx}
\usepackage{textcomp}
\usepackage{xcolor}
\usepackage{float}
\usepackage[utf8]{inputenc}
\usepackage[T1]{fontenc}
\usepackage[portuguese]{babel}

\def\BibTeX{{\rm B\kern-.05em{\sc i\kern-.025em b}\kern-.08em
    T\kern-.1667em\lower.7ex\hbox{E}\kern-.125emX}}
\begin{document}

\title{Binarização Híbrida de Imagens de Documentos Históricos Degradados: Uma Releitura Adaptada\\
{\footnotesize }
}

\author{\IEEEauthorblockN{Vitor Guedes Frade - 221017130 \\ Ryan Reis Fontenele - 211036132 \\ Mateus Elias de Macedo - 222011561}
\IEEEauthorblockA{\textit{Departamento de Ciências da Computação} \\
\textit{Universidade}\\
Brasilia, Brasil \\
vitorguedesfrade@gmail.com \\ 
ryanreisfonte@gmail.com  
\\ mateus.emacedo@gmail.com
}
}

\maketitle

\selectlanguage{english}
\begin{abstract}
This paper presents a reinterpretation and adaptation of a hybrid system for binarizing degraded historical document images. Based on techniques established in the literature, the implemented method integrates pre-processing, dynamic global binarization, adaptive local binarization, and morphological post-processing. Adaptations were made to the original method to improve results, substituting some approaches with optimized implementations, such as Multi-Otsu and Hit-Miss morphology. Parameter adjustments diverging from the base article were empirically chosen to better suit the specific degradation profiles of our test set. The results demonstrate the algorithm's efficacy in preserving writing strokes in low-contrast regions and mitigating background noise, smear, and unwanted artifacts.
\end{abstract}

\begin{IEEEkeywords}
Image Processing, Hybrid Binarization, Historical Documents, Otsu, Nick's Method, CLAHE, Mathematical Morphology.
\end{IEEEkeywords}
\selectlanguage{portuguese}

\section{Introdução}
A binarização de imagens de documentos é uma etapa importante e crítica em muitos sistemas de análise e reconhecimento de documentos, devido a  sua importância para transcrever imagens em níveis de cinza em representações binárias que facilitem tarefas de sistemas de reconhecimento óptico de caracteres (OCR), indexação e preservação digital de documentos. Em documentos históricos, esse processamento é dificultado por degradações típicas: iluminação não uniforme, manchas, vazamento de tinta, variações de contraste e caracteres parcialmente apagados; o que compromete a legilibilidade de algorítmos de reconhecimento, tanto quanto compromete a segmentação do texto. Esta binarização é salutar, devido sua influência sobre todo o processo de segmentação - influenciando a extração de características, leitura e preservação - erros que permaneçam na binarização se propagam na segmentação.

Na literatura de reconhecimento óptico de caracteres, existem abordagens consideradas globais e propostas híbridas que combinam pré-processamento, limiarização e pós-processamento para lidar com regiões da imagem com fatores que dificultariam o reconhecimento óptico pelos algoritmos. Observe-se que o uso de uma abordagem híbrida favorece satisfazer determinados parâmetros que satisfaçam maior legibilidade de acordo com determinados parâmetros, a citar, o contraste.

O presente trabalho consiste na reprodução da metodologia proposta por Boudraa et al. (2019) \cite{boudraa2019}, que sugere um sistema de binarização em múltiplas fases (pré-processamento, binarização híbrida seguida de pós-processamento). Contudo, para fins de melhoria de desempenho e adequação das ferramentas disponíveis, foram implementadas modificações estratégicas na lógica de limiarização global e nas operações de pós-processamento, buscando resultados superiores no tratamento de ruídos e artefatos específicos das imagens utilizadas no conjunto de testes. 

\section{Metodologia e Adaptações}
O desenvolvimento foi conduzido em ambiente Python. O algoritmo foi estruturado nas seguintes etapas, com divergências metodológicas e de parametrização em relação ao artigo base justificadas a seguir:

\subsection{Pré-Processamento e Análise de Contraste}
Para tratar imagens com iluminação deficiente, implementou-se uma etapa de análise de contraste local de Michelson (núcleo $3\times3$). A média global desse mapa de contraste dita a aplicação do CLAHE. O limite de aplicação (contraste médio $< 0,02$) foi mantido como um estimador robusto para imagens altamente desbotadas.

\subsection{Limiarização Global Dinâmica (Divergência TSMO vs. Multi-Otsu)}
Diferentemente do método TSMO original sugerido pela literatura base, esta releitura adota o algoritmo Multi-Otsu (com 3 classes) em conjunto com o Otsu padrão. A decisão sobre qual limiar utilizar baseia-se diretamente no valor do contraste local médio calculado na etapa anterior, utilizando três faixas de restrição empíricas ($T_{01} = 0,03$, $T_{02} = 0,04$ e $T_{03} = 0,085$). 

Esses parâmetros foram ajustados empiricamente para a nossa base de testes, pois os intervalos originais sugeridos no artigo base geravam perdas acentuadas na continuidade dos traços em imagens com contraste marginal. A lógica de decisão estruturada ocorre sob as seguintes faixas:

\begin{itemize}
    \item \textbf{Contraste Extremo Baixo ($\le T_{01}$):} Utiliza o limiar superior do Multi-Otsu.
    \item \textbf{Contraste Muito Baixo ($\le T_{02}$):} Analisa a densidade de pixels na faixa intermediária e a distância entre os limiares. Caso atenda às condições restritivas, adota-se o limiar superior do Multi-Otsu; caso contrário, o Otsu clássico.
    \item \textbf{Contraste Moderado ($\le T_{03}$):} Aplica a limiarização clássica de Otsu.
    \item \textbf{Contraste Alto ($> T_{03}$):} Utiliza o limiar inferior do Multi-Otsu para evitar ruídos de fundo.
\end{itemize}

\subsection{Limiarização Local Adaptativa de Nick}
Para tratar as áreas onde a binarização global falha (textos desbotados e manchas), implementou-se o método adaptativo de Nick:
\begin{equation}
T(x,y) = m + k_{nick} \sqrt{\frac{\sum (p_i^2) - m^2}{N}}
\end{equation}
\textbf{Justificativa do Parâmetro:} O artigo base e a literatura frequentemente sugerem valores de $k_{nick}$ entre $-0,1$ e $-0,2$. Neste trabalho, optou-se pelo valor divergente $k_{nick} = -0,05$ e janela $ws = 7$. Um valor absoluto menor de $k$ torna o limiar menos agressivo. Isso foi intencional: visou-se preservar mais informações dos traços do texto em áreas severamente manchadas, assumindo o risco de manter algum ruído residual, o qual é agressivamente tratado e eliminado na nossa etapa otimizada de pós-processamento (Hit-Miss).

\subsection{Detecção de Manchas e Mesclagem Híbrida}
O algoritmo subdivide a imagem binarizada globalmente em blocos não sobrepostos de $7\times7$. Se a fração de pixels pretos excede $\mu + k_{smear}\sigma$, a região é classificada como mancha. O valor escolhido de $k_{smear} = 0,001$ difere do habitual estatístico (que seria mais conservador) para aumentar a sensibilidade da detecção de manchas (smears), forçando o algoritmo a usar o método de Nick em praticamente qualquer região suspeita de acúmulo de tinta.

\subsection{Pós-processamento Morfológico Hit-Miss e Filtragem de Componentes}
O artigo base sugere alguns passos vagos para o pós processamento da imagem, mas não apresenta muitos detalhes dos métodos específicos utilizados para implementar as soluções. Nessa releitura, foram utilizadas operações morfológicas via transformada \textit{Hit-Miss} para remover buracos e pixels isolados, assim como corrigir leves artefatos de concavidade e convexidade produzidos pela binarização. aliada a uma rigorosa filtragem de componentes conexos. 

Após isso, foi realizada a detecção dos componentes conexos e encontrado suas áreas, como feito no trabalho original, mas algumas mudanças tiveram de ser feitas para resultados mais coerentes serem obtidos. As etapas exatas de toda a pipeline de pós processamento são registrados abaixo:

\begin{enumerate}
    \item \textbf{Gaps e Dots (Pontos isolados e Buracos):} Remoção de falsos pixels de \textit{foreground} e preenchimento de buracos de 1 pixel no texto utilizando matrizes específicas da transformada \textit{Hit-Miss}.
    \item \textbf{Filtragem por Componentes Conexos:} O algoritmo utilizado analisa todas as componentes e remove sumariamente aquelas com áreas inferiores a um limite estatístico restritivo $(\lambda \cdot \mu) / \sigma$, onde $\mu$ é a média do tamanho das áreas dos componentes, $\sigma$ seu desvio padrão e $\lambda = 15$ uma constante empírica. No trabalho original, as equações propostas implicavam na remoção dos componentes de área maior que a expressão $(\lambda \cdot \mu) / \sigma$, porém a implementação através dessa lógica resultava na remoção quase inteira do texto em especial em casos de escrita cursiva, nos quais palavras inteiras formam um componente único. Devido à escassez de detalhes dessa etapa no artigo e à ausência de um repositório acessível, não se pode confirmar se essa lógica é intencional, ou resultado de um erro de transcrição, mas sua inversão resultou em notável aprimoramento da imagem final e, portanto, escolheu-se eliminar os componentes menores.
    
    \item \textbf{Convexidades e Concavidades:} Por fim, foram criados bancos de kernels que representavam concavidades e convexidades de um pixel em cada direção (Cima, Baixo, Esquerda, Direita). Essas matrizes foram utilizadas pela operação Hit-Miss suavizar o contorno dos caracteres e remover artefatos gerados pela binarização.
\end{enumerate}

\subsection{Métodos de Avaliação}

Para avaliar o desempenho do algoritmo em sua binarização, foram  utilizadas quatro métricas principais, sendo feitas com comparação a\textit{ground truths} disponibilizadas junto às bases de dados utilizadas. As métricas utilizadas foram:

\begin{enumerate}
    \item \textbf{\textit{F-Measure}:} Uma medição padrão da acuracia de módelos representada pela fórmula:
    
    \[FM = \frac{2*Recall*Precision}{Recall+Precision}\] 

    onde $Recall = \frac{TP}{TP+FN}$ e $Precision = \frac{TP}{TP+FP}$, sendo $TP$,$FN$ e $FP$ respectivamente os valores de verdadeiro positivo, falso negativo e falso positivo do modelo.

    \item \textbf{\textit{pseudo F-Measure}:} Uma medida introduzida em \cite{pfm} e baseada o \textit{F-Measure}, mas que substitui Recall por pseudo-Recall $R_{ps}$ e Precision por pseudo-Precision $P_{ps}$. Essas medidas são baseadas na distância ao contorno dos caracteres da \textit{ground truth} e, no caso do pseudo-Recall, também é normalizado de acordo com a espessura do traço.

    \item \textbf{\textit{Peak Signal to Noise Ratio (PSNR)}:} Uma medida que determina a semelhança entre duas imagens,  sendo $I$ e $I'$ duas imagens de dimensão $M$×$N$ e $C$ a diferença entre pixels do \textit{foreground} e \textit{background}, ela é definida como:

    \[PSNR = 10log(\frac{C^2}{MSE})\]

    onde \[MSE = \frac{\sum_{x=1}^{M}\sum_{y}^{n}(I(x,y) - I'(x,y)^2}{MN}\]
    \item \textbf{\textit{Distance Reciprocal Distortion Metric (DRD)}:} Uma medida baseada na visão humana que mede a distorção visual em documentos, como usada em \cite{drd}. A distorção de todos os $S$ \textit{flipped pixels} é definida por:

    \[DRD = \frac{\sum_{k=1}^{S}DRD_{k}}{NUBN}\]

    Nessa expressão, $NUBN$ é o número de blocos $8$×$8$ não uniformes na imagem e $DRD_{k}$ é a distorção do k-ésimo pixel virado, calculado através da matriz normalizada $5$×$5$ de pesos $W_{NM}$, definida em \cite{drd}.

    
\end{enumerate}

\section{Resultados e Discussão}
O método proposto foi testado em algumas imagens mais simples de teste, mas também na base de imagens H-DIBCO de 2014 \cite{hdibco2014}, composta por 10 imagens de documentos manuscritos, suas \textit{ground truths} e executáveis utilizados para avaliar o desempenho do algoritmo testado através das métricas previamente apresentadas. 

À medida que os resultados foram considerados, foi possível concluir que, mesmo não sendo uma métrica principal, certos métodos apresentaram desempenho computacional melhor que outros sem apresentar perda em acurácia, o que auxiliou no desenvolvimento. Em um desses casos, por exemplo, o cálculo de contraste local aliado à decisão condicional entre Otsu e Multi-Otsu mostrou-se computacionalmente mais leve e eficiente que o TSMO original para as imagens do conjunto de nosso conjunto de dados.

Quanto à qualidade em si da binarização produzida, os benefícios dos métodos utilizados comparados aos tradicionais foram exemplificadas por certas imagens em condições mais precárias. Por exemplo, a Figura \ref{fig:h05} ilustra perfeitamente a necessidade da binarização híbrida e o papel crucial da binarização de Nick nas manchas (smear). Na binarização puramente global (Fig. \ref{fig:h05}b), grandes áreas de acúmulo de tinta ou variação de iluminação colapsam em manchas densas, ocultando completamente o texto. Ao identificar essas regiões (densidade $> \mu + k_{smear}\sigma$) e aplicar localmente a fórmula de Nick com nosso $k_{nick}=-0,05$ otimizado (Fig. \ref{fig:h05}c), o texto é resgatado do fundo escuro. Por fim, o pós-processamento Hit-Miss (Fig. \ref{fig:h05}d) elimina os resíduos pontuais e preserva a integridade estrutural da escrita. Isso pode ser visto nas Figuras 1-4.

\begin{figure}[htbp]
    \centering
    \begin{minipage}{0.48\columnwidth}
        \centering
        \includegraphics[width=\linewidth]{resultados_binarizacao_nick/H05_01_Original_Gray.png}
        \vspace{-0.2cm}
        \caption*{(a) Imagem Original}
    \end{minipage}
    \hfill
    \begin{minipage}{0.48\columnwidth}
        \centering
        \includegraphics[width=\linewidth]{resultados_binarizacao_nick/H05_02_Binarizada_global.png}
        \vspace{-0.2cm}
        \caption*{(b) Binarização Global}
    \end{minipage}
    \vspace{0.3cm}
    
    \begin{minipage}{0.48\columnwidth}
        \centering
        \includegraphics[width=\linewidth]{resultados_binarizacao_nick/H05_03_Binarizada_Nick.png}
        \vspace{-0.2cm}
        \caption*{(c) Binarização de Nick}
    \end{minipage}
    \hfill
    \begin{minipage}{0.48\columnwidth}
        \centering
        \includegraphics[width=\linewidth]{resultados_binarizacao_nick/H05_04_Pos_Processada.png}
        \vspace{-0.2cm}
        \caption*{(d) Final (Pós-processada)}
    \end{minipage}
    \caption{Comparativo das etapas na imagem H05. A binarização global (b) deixa manchas opacas devido à degradação. A aplicação da limiarização de Nick (c) nestas regiões recupera o texto, sendo o ruído residual limpo no pós-processamento (d).}
    \label{fig:h05}
\end{figure}

Caso o algoritmo fosse executado com os parâmetros sugeridos pelo artigo base (por exemplo, valores de $k_{nick}$ mais negativos e janelas maiores), notaríamos uma erosão excessiva dos caracteres resgatados nas áreas de mancha. A nossa parametrização compensa isso balanceando um $k_{nick}$ tolerante com uma limpeza morfológica posterior rigorosa. 

Entretanto, foram percebidos alguns desafios quando aplicado nosso algoritmo sobre as imagens de H-DIBCO 2014 \cite{hdibco2014}. 

Para a medição do desempenho, foram avaliadas as métricas para o processo para ambas a binarização padrão e a Nick. Os resultados desse processo para a binarização global e Nick, assim como os resultados consolidades de vários métodos apresentados no artigo original (Entre os quais \textit{Boudra} é o proposto pelos autores) são representados nas Figuras 6, 7 e 8 respectivamente.

\begin{figure}[htbp]
    \centering
    \begin{minipage}{0.9\columnwidth}
        \centering
        \includegraphics[width=\linewidth]{Avaliacao/Tabela_Global.png}
        \vspace{-0.2cm}
        \caption*{Tabela das métricas da binarização sem método Nick}
    \end{minipage}
\end{figure}

\begin{figure}[htbp]
    \centering
    \begin{minipage}{0.9\columnwidth}
        \centering
        \includegraphics[width=\linewidth]{Avaliacao/Tabela_Nick.png}
        \vspace{-0.2cm}
        \caption*Tabela das métricas da binarização com Nick
    \end{minipage}
\end{figure}

\begin{figure}[htbp]
    \centering
    \begin{minipage}{0.9\columnwidth}
        \centering
        \includegraphics[width=\linewidth]{Avaliacao/Tabela_Artigo.png}
        \vspace{-0.2cm}
        \caption*Tabela das métricas apresentadas no artigo original.
    \end{minipage}
\end{figure}

Através desses resultados, é possível ver alguns pontos notáveis de divergência entre as metodologias utilizadas. Apesar da importância da abordagem Nick para lidar com manchas e outras anomalias nos documentos, nossa implementação utilizada possui, em média, um desempenho significativamente inferior para a base de dados H-DIBCO 2014.

Esse padrão é alarmante, mas é importante compreender a raiz dessa discrepância, o que pode ser feito através da análise dos caso mais excepcionais. A imagem 6 e suas binarizações são visíveis nas figuras 9-11 e servem para demonstrar a falha fatal da implementação do método Nick realizada.

\begin{figure}[htbp]
    \centering
    \begin{minipage}{0.8\columnwidth}
        \centering
        \includegraphics[width=\linewidth]{resultado_HDIBCO/H06_01_Original_Gray.png}
        \vspace{-0.2cm}
        \caption*{(a) Imagem Original}
    \end{minipage}
    \hfill
    \begin{minipage}{0.8\columnwidth}
        \centering
        \includegraphics[width=\linewidth]
        {resultado_HDIBCO/H06_GPos_processada.png}
        \vspace{-0.2cm}
        \caption*{(b) Binarização Global}
    \end{minipage}
    \vspace{0.3cm}
    
    \begin{minipage}{0.8\columnwidth}
        \centering
        \includegraphics[width=\linewidth]{resultado_HDIBCO/H06_NPos_processada.png}
        \vspace{-0.2cm}
        \caption*{(c) Binarização de Nick}
    \end{minipage}
    \caption{Comparativo da imagem e suas binarizações}
    \label{fig:h06-bin}
\end{figure}

Nesse caso, a \textit{F-measure} do métoo Nick foi a extremamente baixa 38.7276, comparada ao valor de 93.7172 da binarização global. Apesar desse caso ser extremo, ele demonstra que nossa implementação Nick sofre com escrita de traço mais espesso, preservando apenas seu contorno e não seu interior, o que afeta suas métricas significativamente.

Apesar da causa exata dessa discrepância não ter sido encontrada diretamente, modificações nos parâmetros das operações, assim como usos de métodos de preenchimento poderiam aprimorar seu resultado e fazer sua métricas alcançarem ou ultrapassarem suas alternativas.

Quanto à binarização globalizada, é possível notar que, para essa base de dados, seus resultados são incrivelmente professores. Apesar de não ultrapassar o método de \textit{Boudra} proposto no artigo, seu desempenho alcança e, em algumas métricas, supera o método \textit{Otsu}, o segundo melhor. Isso demonstra que foi possível replicar vários dos fatores que tornaram o algoritmo do artigo original eficiente, mesmo que certos aspectos não terem sido possíveis de recriar com precisão completa.

\section{Conclusão}
A partir da releitura do artigo, o intuito inicial foi obter resultados semelhantes aos apresentados Boudraa et al. (2019) \cite{boudraa2019}, ao qual inicialmente seguimos escrupulosamente sua metodologia. Ao realizarmos os testes com o banco de imagens citado (dataset), observou-se que alguns parâmetros e, principalmente a etapa de pós-processamento, exigiam ajustes para melhor adequação às características específicas das imagens avaliadas.
Diante disto, foram investigados parâmetros e técnicas que favorecessem a legibilidade do texto e a preservação das informações relevantes para um futuro reconhecimento óptico de caracteres (OCR), processo este no qual se verificou que o algoritmo de Nick, apesar de eficiente na delimitação de traços de contorno, apresentou limitações na região do interior dos contornos, sobretudo em regiões de escrita mais espessa. Apesar deste fator, o método global mostrou bom desempenho e, em situações com manchas e desgastes acentuados, o método adaptativo de Nick desempenhou nestas situações potencial satisfatório, apesar de demandar correções.

A adaptação consciente dos parâmetros, sobretudo no limiar do algoritmo de Nick e no critério estatístico de detecção de manchas, somada à substituição do pós-processamento original por operações Hit-Miss, demonstrou-se como solução eficaz para os testes realizados. Os resultados obtidos indicam recuperação satisfatória de informações em documentos gravemente degradados, evidencias da validade técnica das adaptações realizadas. 

Aos trabalhos futuros, recomenda-se a reimplementação do algoritmo de Nick, de modo a sanar a preservação do interior dos traços em regiões com escrita mais espessa, além da exploração de técnicas de aprendizado profundo para aprimorar a detecção de manchas e o processo de binarização adaptativa.
\begin{thebibliography}{00}
\bibitem{boudraa2019}
O. Boudraa, W. K. Hidouci, and D. Michelucci, ``Degraded Historical Documents Images Binarization Using a Combination of Enhanced Techniques,'' \textit{arXiv preprint arXiv:1901.09425}, 2019.
\bibitem{hdibco2014}
K. Ntirogiannis, B. Gatos and I. Pratikakis, "ICFHR2014 Competition on Handwritten Document Image Binarization (H-DIBCO 2014)," 2014 14th International Conference on Frontiers in Handwriting Recognition, Hersonissos, Greece, 2014, pp. 809-813, doi: 10.1109/ICFHR.2014.141.
\bibitem{pfm}
Chris Tensmeyer and Tony Martinez. 2020. Historical Document Image Binarization: A Review. SN Comput. Sci. 1, 3 (May 2020). https://doi.org/10.1007/s42979-020-00176-1
\bibitem{drd}
Lu, Haiping \& Kot, A.C. \& Shi, Y.Q.. (2004). Distance-Reciprocal Distortion Measure for Binary Document Images. IEEE Signal Processing Letters. 11. 228-231. 10.1109/LSP.2003.821748. 
\end{thebibliography}

\end{document}
"""

with open("relatorio_releitura_v2.tex", "w", encoding="utf-8") as f:
    f.write(latex_content)

print("Fetched content: Successfully generated relatorio_releitura_v2.tex")
print("[file-tag: relatorio_releitura_v2.tex]")